\section{Conclusion}
\label{sec:conclusion}

This paper documents that board observer networks serve as channels for private information flow between firms. When material events occur at private companies with board observers, connected public companies show abnormal returns of 4.6 percentage points (M\&A announcements) to 9.5 percentage points (bankruptcy filings) in the weeks before public disclosure. The effect is concentrated in same-industry connections---where the information is most relevant---and in events involving confidential board decisions. Routine corporate events produce no spillover.

Two regulatory changes provide causal evidence. The NVCA's 2020 removal of fiduciary language from observer provisions coincides with an increase in same-industry spillover from zero to 3.0 percentage points. The DOJ/FTC's January 2025 extension of Clayton Act Section~8 to observers coincides with a reversal. These opposite-direction responses to independent regulatory changes, both with clean placebos, support a causal interpretation: the accountability structure governing observers---not unobserved firm characteristics or market trends---shapes the information channel.

Our findings carry direct regulatory implications. The DOJ/FTC's decision to extend interlocking directorate rules to board observers is empirically supported: same-industry observer connections generate significant abnormal returns at connected firms months before public disclosure. The NVCA's 2025 expansion of competitive harm carve-outs, which gives companies broader authority to exclude observers from sensitive discussions, responds to a real information channel documented in our data.

Several limitations warrant acknowledgment. Our network, while supplemented with BoardEx and Form~4 data, captures an estimated half of all observer-to-public-firm connections; the remaining connections would increase statistical power but cannot introduce spurious effects. The same-industry samples for M\&A and bankruptcy events are small (11--28 connected same-industry observations), making coefficient estimates noisy despite their statistical significance. We do not identify the specific mechanism through which information reaches prices---whether through trading by the observer, strategic decisions at the connected firm, or broader network diffusion---though the regulatory responsiveness suggests the observer's legal position is a necessary condition.

Future research could extend this analysis in several directions: incorporating events at public observed companies (which triples the sample), adopting a calendar-time portfolio approach as an alternative to event-time CARs, and examining whether the October 2025 NVCA competitive harm provisions reduce spillover as additional post-shock data accumulates.